\newcommand{\inputparam}[2]{%
  \def\parg{#2}
  \input{#1}
}
\inputparam{intestazione}{Idea progettuale}
\newpage
\section{Natura dell'applicazione}
Evnt è un’applicazione mobile e web per la scoperta, la creazione e la gestione di eventi locali. L’obiettivo è permettere agli utenti di trovare attività vicine ai propri interessi e alla propria posizione, partecipare agli eventi, salvarli tra i preferiti, comunicare con gli altri partecipanti e creare nuovi eventi in modo semplice. L’app integra feed, mappa interattiva, creazione guidata degli eventi, chat e notifiche in-app.
\section{Il problema che risolve}
L'app risponde alla crescente frammentazione dell’informazione digitale e al fenomeno dell'isolamento sociale paradossale. Da un lato, risolve l'inefficienza dei metodi di ricerca tradizionali (spesso dispersivi e non profilati); dall’altro, combatte la difficoltà di trovare community affini. Attraverso un sistema di ricerca granulare e filtri basati sugli interessi, l'app aggrega gli individui non solo attorno a un luogo, ma attorno a una preferenza condivisa, facilitando una socializzazione spontanea e di qualità.

\section{Il target di utenti}
L’applicazione è rivolta al target 16+, cercando di caratterizzare gli eventi in base all’età.

\section{Le macrofunzionalità previste}
Il sistema si articola in cinque macrofunzionalità principali, pensate per coprire l’intero ciclo di vita dell’esperienza utente: dalla scoperta degli eventi alla partecipazione, fino alla comunicazione e alla gestione del proprio profilo.

\begin{enumerate}
    \item \textbf{Scoperta personalizzata degli eventi} \\
    L’app consente agli utenti di trovare eventi pertinenti in base ai propri interessi, alla posizione geografica o alla città dichiarata in fase di registrazione. La ricerca può essere raffinata tramite filtri su categorie, prezzo, distanza e testo libero.

    \item \textbf{Contestualizzazione geografica} \\
    Gli eventi vengono associati a un luogo reale e visualizzati in base alla vicinanza rispetto all’utente. In caso di mancato consenso alla geolocalizzazione, il sistema utilizza la città indicata nel profilo come riferimento alternativo.

    \item \textbf{Creazione e gestione degli eventi} \\
    Gli utenti possono creare eventi specificando titolo, descrizione, categoria, sottocategoria, indirizzo, data, orario, numero massimo di partecipanti, costo e modalità di comunicazione. L’organizzatore può modificare o annullare gli eventi creati.

    \item \textbf{Partecipazione e salvataggio} \\
    L’utente può consultare i dettagli di un evento, iscriversi, rimuovere la propria iscrizione e salvarlo tra i preferiti. Il sistema mantiene aggiornate le informazioni relative a partecipanti, disponibilità dei posti e stato dell’evento.

    \item \textbf{Comunicazione e notifiche interne} \\
    La piattaforma supporta la comunicazione tra utenti tramite chat collegate agli eventi e conversazioni private avviate tramite email. Le notifiche in-app informano l’utente su aggiornamenti, cancellazioni, messaggi e promemoria rilevanti.
\end{enumerate}

\section{L'analisi dei competitor}
\begin{table}[h]
    \centering
    \begin{tabular}{|p{3cm}|p{5cm}|p{5cm}|} 
        \hline
        \textbf{Nome} & \textbf{Punti di forza} & \textbf{Limiti} \\
        \hline
        
        Eventbrite & 
        \begin{itemize}[nosep, leftmargin=*]
            \item È la più famosa
            \item È sicura per comprare biglietti di ogni tipo
        \end{itemize} & 
        \begin{itemize}[nosep, leftmargin=*]
            \item È un po' fredda e troppo professionale
            \item Le commissioni sui biglietti sono alte
        \end{itemize} \\ \hline
        
        Soiree & 
        \begin{itemize}[nosep, leftmargin=*]
            \item Ottima per conoscere persone nuove in modo naturale e partecipare a feste curate
        \end{itemize} & 
        \begin{itemize}[nosep, leftmargin=*]
            \item C'è meno scelta rispetto ai colossi
            \item Funziona bene solo nelle grandi città
        \end{itemize} \\ \hline
      
        Meetup & 
        \begin{itemize}[nosep, leftmargin=*]
            \item Ideale per trovare gruppi di persone con stessi hobby
        \end{itemize} & 
        \begin{itemize}[nosep, leftmargin=*]
            \item App graficamente vecchia
            \item Gruppi spesso inattivi
        \end{itemize} \\ \hline

        Fever &
        \begin{itemize}[nosep, leftmargin=*]
            \item Ti suggerisce le cose e gli eventi di tendenza in città
        \end{itemize} &
        \begin{itemize}[nosep, leftmargin=*]
            \item Spinge molto i propri eventi e mostra solo ciò che è di moda, non tutto quello presente
        \end{itemize} \\ \hline
        
    \end{tabular}
\end{table}

\end{document}

